\section{Discussion\label{chapter5}}

This chapter interprets the experimental results, analyzes the key trade-offs, and positions the findings within the broader literature. Section~\ref{sec:key_findings} discusses the dataset-dependent performance pattern. Section~\ref{sec:budget_discussion} examines the parameter budget matching methodology. Section~\ref{sec:monotonic_trends} analyzes the observed monotonic hyperparameter trends in ESNs on chaotic time series. Section~\ref{sec:efficiency_discussion} analyzes training efficiency trade-offs. Section~\ref{sec:limitations} acknowledges the limitations of this study. Section~\ref{sec:literature_comparison} compares the results with prior work.

% ============================================================================
\subsection{Interpretation of Key Findings}
\label{sec:key_findings}
% ============================================================================

The most striking result of this thesis is the strongly dataset-dependent nature of model performance. PyReCo (ESN) achieves near-perfect predictions on the Lorenz (R\textsuperscript{2} = 0.999) and Mackey-Glass (R\textsuperscript{2} = 0.987) systems, while LSTM dominates on the Santa Fe laser dataset (R\textsuperscript{2} = 0.934 vs.\ 0.770). This pattern is consistent across all parameter budgets, data lengths, and random seeds, with large effect sizes ($d > 2.0$), ruling out chance variation.

\subsubsection{Why PyReCo Excels on Synthetic Chaotic Systems}

The Lorenz and Mackey-Glass systems share key characteristics that favor ESNs:
\begin{itemize}
    \item \textbf{Deterministic dynamics}: Both systems are generated by known equations with no measurement noise. The reservoir's fixed nonlinear dynamics can ``resonate'' with the attractor structure, creating internal representations that naturally capture the system's geometry.
    \item \textbf{Low-dimensional attractors}: The Lorenz attractor is 3-dimensional, and the Mackey-Glass attractor, while technically infinite-dimensional due to the delay, has a relatively low fractal dimension ($d_f \approx 2.1$ for $\tau = 17$). The reservoir, even with modest size, provides sufficient dimensionality to represent these attractors.
    \item \textbf{Smooth dynamics}: Both systems exhibit smooth, continuous trajectories without discontinuities or rapid transients. The reservoir's leaky-integrator dynamics (Eq.~\ref{eq:esn_state}) are well-suited to tracking such smooth evolution.
\end{itemize}

The ridge regression readout then performs a simple linear mapping from the reservoir state to the prediction, which is sufficient when the reservoir dynamics provide a good basis for the target signal.

\subsubsection{Why LSTM Excels on Real-World Data}

The Santa Fe laser dataset differs fundamentally from the synthetic systems:
\begin{itemize}
    \item \textbf{Measurement noise}: Experimental data inevitably contains sensor noise and quantization effects not present in simulations. LSTM's trainable parameters can learn to filter noise through gradient-based optimization, while the fixed reservoir has no such adaptation mechanism.
    \item \textbf{Unknown and complex dynamics}: The NH\textsubscript{3} laser system has a high-dimensional state space. The reservoir's random projection may not capture the relevant features as effectively as LSTM's learned representations.
    \item \textbf{Non-stationary patterns}: The Santa Fe data exhibits collapse events with irregular timing. LSTM's gating mechanism enables it to learn complex conditional dynamics (e.g., ``collapse is imminent when these features align''), while the ESN's fixed reservoir cannot adapt its dynamics to such patterns.
\end{itemize}

This interpretation aligns with the broader machine learning principle that learned representations generally outperform fixed random projections when the data distribution is complex and noisy, while fixed projections can be more efficient when the data lies on a low-dimensional manifold with clean structure.

Importantly, even on the Lorenz system where PyReCo normally excels, the ESN's advantage breaks down at very low training fractions. The supplementary broad grid search (Section~\ref{sec:data_eff}) confirms that this degradation is an architectural limitation rather than a hyperparameter mismatch: expanding the search from 6 to 180 hyperparameter combinations at train\_frac = 0.2 yields only marginal improvement (R\textsuperscript{2} $<$ 0.5 in all cases). The fixed reservoir requires a minimum amount of data to adequately cover the attractor, below which even optimal hyperparameters cannot compensate.

% ============================================================================
\subsection{Parameter Budget Analysis}
\label{sec:budget_discussion}
% ============================================================================

A central methodological contribution of this thesis is matching models by \emph{total} parameter count rather than trainable parameter count. This choice has important implications.

At the small budget (${\sim}$1K parameters), the ESN has 100 trainable readout weights while the LSTM has 952 fully trainable parameters---a 9.5$\times$ difference in optimization degrees of freedom. Despite this, the ESN achieves higher R\textsuperscript{2} on Lorenz and Mackey-Glass. This demonstrates that for chaotic systems with low-dimensional attractors, the reservoir's fixed nonlinear expansion is more valuable than additional trainable parameters.

At the large budget (${\sim}$50K parameters), the ESN has 700 trainable weights while the LSTM has 49\,077---a 70$\times$ difference. On the Santa Fe dataset, LSTM's advantage grows with budget size, suggesting that the additional trainable parameters enable it to learn increasingly refined representations of the complex real-world signal.

An alternative matching strategy---equalizing \emph{trainable} parameters---would require either a very large reservoir (to give 49\,077 readout weights) or a very small LSTM (to match 700 trainable parameters). Both alternatives would introduce different biases. The total parameter budget was chosen because it equalizes the overall model footprint (memory, matrix dimensions), which is arguably the fairest comparison of model capacity \cite{wringe2025benchmarks}.

% ============================================================================
\subsection{Monotonic Hyperparameter Trends on Chaotic Time Series}
\label{sec:monotonic_trends}
% ============================================================================

An important observation from the hyperparameter search is that, across all three datasets, the four ESN hyperparameters exhibit \emph{monotonic} effects on prediction performance within the tested ranges. This explains why the optimal hyperparameter configurations consistently fall at the boundaries of the search grid.

Table~\ref{tab:monotonic} summarizes the observed trends and their theoretical basis.

\begin{table}[H]
\centering
\caption{Monotonic hyperparameter trends observed on chaotic time series.}
\label{tab:monotonic}
\small
\begin{tabular}{@{}llll@{}}
\toprule
Parameter & Observed Trend & Optimal Direction & Literature Support \\
\midrule
Spectral radius     & Monotonically increasing & $\to 0.99$ & Edge of chaos theory \cite{bertschinger2004real} \\
Leakage rate        & Monotonically increasing & $\to 0.7$--$0.8$ & Timescale matching \cite{lukosevicius2012practical} \\
Density             & Monotonically decreasing & $\to 0.01$--$0.03$ & Implicit regularization \cite{lukosevicius2009reservoir} \\
Input fraction      & Monotonically decreasing & $\to 0.1$--$0.3$ & Input scaling theory \cite{lukosevicius2012practical} \\
\bottomrule
\end{tabular}
\end{table}

\textbf{Spectral radius $\to$ near 1.0.} The edge of chaos theory predicts that ESNs perform best when the spectral radius approaches unity, where the reservoir maintains rich, long-lived dynamics without becoming unstable \cite{bertschinger2004real}. Chaotic time series with long-range dependencies benefit from the extended memory capacity that comes with high spectral radius. All three datasets consistently selected the highest available spectral radius value.

\textbf{Leakage rate $\to$ high values.} The optimal leakage rate depends on the timescale of the input signal relative to the reservoir update rate. For chaotic systems sampled at typical rates, higher leakage (faster state updates) is preferred because the reservoir must track rapidly evolving dynamics \cite{lukosevicius2012practical}.

\textbf{Density $\to$ sparse.} Sparse connectivity provides implicit regularization and promotes diverse dynamic substructures within the reservoir. In practice, very sparse reservoirs (1--5\% connectivity) often perform comparably to or better than dense ones \cite{lukosevicius2009reservoir}, and our experiments confirm this for chaotic prediction tasks.

\textbf{Input fraction $\to$ low values.} Reducing input connectivity forces the reservoir to develop richer internal representations through recurrent dynamics rather than directly propagating input signals. This is consistent with input scaling analysis in Luko\v{s}evi\v{c}ius \cite{lukosevicius2012practical}.

\textbf{Dataset-dependent parameter sensitivity.} While the monotonic \emph{directions} are consistent across all three datasets, the \emph{magnitude} of sensitivity varies substantially, as summarized in Table~\ref{tab:sensitivity}.

\begin{table}[H]
\centering
\caption{Dataset-dependent ESN hyperparameter sensitivity.}
\label{tab:sensitivity}
\small
\begin{tabular}{@{}llll@{}}
\toprule
Dataset & $\Delta R^2$ Across Grid & Sensitivity & Key Characteristics \\
\midrule
Lorenz       & $< 0.01$     & Low      & Deterministic 3D ODE, noise-free \\
Mackey-Glass & $0.05$--$0.15$ & Moderate & DDE ($\tau = 17$), infinite-dimensional \\
Santa Fe     & $0.10$--$0.30$ & High     & Real experimental data, noisy, small \\
\bottomrule
\end{tabular}
\end{table}

The Lorenz system, despite being chaotic, has a highly regular attractor structure (butterfly-shaped, 3-dimensional). Noise-free numerical simulation with high signal-to-noise ratio means that even suboptimal reservoir configurations can capture this regular structure, making performance robust to hyperparameter choices. The Mackey-Glass equation, a delay differential equation with $\tau = 17$, is effectively infinite-dimensional, with longer and more complex temporal dependencies. The reservoir must precisely match the delay timescale---the leakage rate and spectral radius directly determine whether the reservoir can retain information from $\tau$ steps ago \cite{jaeger2004harnessing}. The Santa Fe laser dataset, consisting of real experimental measurements, contains measurement noise and unknown high-dimensional dynamics. The limited dataset size (${\sim}$10\,000 points) amplifies overfitting risk: slight parameter deviations can cause the reservoir to either overfit noise or underfit the signal \cite{weigend1994time}.

The general principle is that ESN parameter sensitivity correlates with data complexity, noise level, and dataset size. This gradient---from pure numerical simulation (Lorenz) through delay-based simulation (Mackey-Glass) to real experimental data (Santa Fe)---mirrors the gradient in prediction performance (Section~\ref{sec:key_findings}), suggesting that the same data characteristics that make a problem harder for ESNs also make them harder to tune.

\textbf{Implications for hyperparameter search.} When a parameter's effect on performance is monotonic within the search range, the optimum will always appear at the grid boundary. Extending the grid shifts the boundary but cannot change the monotonic nature---only physical or mathematical constraints (spectral radius $< 1.0$ for guaranteed echo state property, density $> 0$, input fraction $> 0$) define the true limits. This observation has two practical consequences: (1) narrower, boundary-focused search grids could reduce tuning cost without sacrificing performance on chaotic tasks; and (2) the monotonic trends are specific to \emph{chaotic time series with long-range dependencies}---for other task types (e.g., slowly varying signals, classification), the optimal directions may differ. Furthermore, the dataset-dependent sensitivity suggests that the tuning budget should be allocated proportionally to data complexity: low-sensitivity systems like Lorenz tolerate coarse grids, while high-sensitivity real-world data like Santa Fe demands finer resolution or adaptive search strategies.

% ============================================================================
\subsection{Training Efficiency Trade-offs}
\label{sec:efficiency_discussion}
% ============================================================================

By separating training time into three components---training with optimal parameters, hyperparameter tuning, and inference---the efficiency comparison becomes more nuanced than a single aggregate number would suggest.

\textbf{Training with optimal parameters (primary metric)}: The crossover pattern is clear: PyReCo is 32$\times$ faster at small scale but LSTM is 2.4$\times$ faster at large scale. At small budgets, ridge regression on a 100-node reservoir completes in 1.7 seconds compared to 56 seconds for LSTM's gradient-based training. At large budgets, the $O(N^2)$ cost of a 700-node reservoir (74.7 seconds) exceeds LSTM training (30.6 seconds). This crossover reflects a fundamental algorithmic trade-off: ridge regression scales quadratically with reservoir size, while LSTM training scales more favorably with hidden size due to efficient GPU-optimized matrix operations.

\textbf{Hyperparameter tuning (supplementary)}: The tuning cost reveals an important asymmetry. PyReCo requires a larger search grid (36 combinations with 5-fold CV) than LSTM (9 combinations with a single validation split), reflecting the documented higher sensitivity of ESNs to hyperparameter choices \cite{matzner2022hyperparameter}. At the large budget, PyReCo's tuning cost (43 minutes) dominates the total computational budget, while LSTM's tuning remains modest (5 minutes). However, this cost is incurred once per configuration and is amortized in deployment scenarios where models are retrained with fixed hyperparameters.

\textbf{Inference time}: LSTM is 50--560$\times$ faster at inference across all budgets. This advantage arises because LSTM inference involves a compact hidden state ($h = 10$--$109$), while ESN inference requires updating a much larger reservoir state ($N = 100$--$700$). For latency-critical or throughput-intensive applications, this difference is decisive.

\textbf{Training efficiency (R\textsuperscript{2}/second)}: Using training time with optimal parameters as the denominator, PyReCo achieves 22$\times$ more R\textsuperscript{2} per second at small scale. This metric captures the ``return on computational investment'' and strongly favors ESNs for resource-constrained applications where rapid retraining is needed.

% ============================================================================
\subsection{Limitations}
\label{sec:limitations}
% ============================================================================

Several limitations should be considered when interpreting these results:

\begin{enumerate}
    \item \textbf{Single-layer LSTM architecture}: Only single-layer (and implicitly two-layer) LSTM architectures were evaluated. Modern sequence models including multi-layer LSTMs, GRUs, attention mechanisms, and Transformers might yield different results, particularly on the Santa Fe dataset.

    \item \textbf{Univariate prediction}: Only scalar (univariate) prediction was evaluated. For the Lorenz system, which is intrinsically 3-dimensional, multivariate prediction could reveal different performance characteristics.

    \item \textbf{Three benchmark datasets}: While the three datasets span synthetic and real-world chaotic systems, generalizing to other domains (weather, finance, biological systems) requires additional validation.

    \item \textbf{pyReCo library performance}: The large-scale training time disadvantage of PyReCo is partly attributable to the library implementation rather than algorithmic limitations. A custom ESN implementation with optimized matrix operations might change the efficiency comparison at large scale.

    \item \textbf{Total vs.\ trainable parameter matching}: The choice to match total parameters rather than trainable parameters affects the comparison. Under trainable parameter matching, ESNs would use much larger reservoirs, potentially improving their performance at the cost of increased memory usage.

    \item \textbf{No streaming/online evaluation}: All experiments use batch training and evaluation. In real-time applications, ESNs could potentially be updated incrementally, while LSTMs require more complex online learning procedures.

    \item \textbf{Hyperparameter search asymmetry}: PyReCo uses 36 hyperparameter combinations with 5-fold cross-validation (180 training runs) while LSTM uses 9 combinations with a single validation split (9 training runs). Although this reflects the documented higher sensitivity of ESNs to hyperparameters \cite{matzner2022hyperparameter}, it introduces an asymmetry in computational effort dedicated to model selection. By reporting tuning cost separately from training time with optimal parameters (Section~\ref{sec:train_eff}), this asymmetry is made transparent rather than hidden in a single aggregate metric.
\end{enumerate}

% ============================================================================
\subsection{Comparison with Literature}
\label{sec:literature_comparison}
% ============================================================================

The findings of this thesis are broadly consistent with prior comparative studies. Chattopadhyay et al. \cite{chattopadhyay2020data} found that RC achieves the best accuracy-to-cost ratio for the Lorenz 96 system, aligning with our finding that PyReCo dominates on the Lorenz system. Vlachas et al. \cite{vlachas2020backpropagation} reported that RC better reproduces the long-term climate of chaotic attractors, which is consistent with PyReCo's superior single-step performance on synthetic data.

The dataset-dependent performance pattern echoes the observation by Jirak et al. \cite{jirak2023echo} that fair comparison between ESN and LSTM is challenging because performance rankings depend on the specific task. Our systematic evaluation across three datasets with matched budgets provides quantitative evidence for this claim.

The key \textbf{novel contribution} of this thesis is the matched total parameter budget methodology. While prior studies have compared models of different sizes or without controlling for model capacity, our approach isolates the effect of architecture from the effect of model size. The finding that ESNs achieve superior performance with only 1--10\% trainable parameters (on synthetic data) provides strong evidence for the computational efficiency of the reservoir computing paradigm for well-characterized dynamical systems.
